% Actualizacion 2024: Logotipos de la UPV


\documentclass[12pt]{article}
\usepackage[left=2.5cm,top=2.0cm,right=2.5cm,bottom=3.0cm]{geometry}
\usepackage[utf8]{inputenc}
\usepackage[spanish]{babel}
%\usepackage[english, american]{babel}
%\usepackage[spanish,es-tabla]{babel}
\usepackage[linguistics]{forest}
\usepackage{amssymb, amsmath, amsbsy} % simbolitos
\usepackage{longtable} % para tablas largas
\usepackage{graphicx}
\usepackage{fancyhdr}
\usepackage{xcolor}
\usepackage{multirow}
\usepackage{listings}
\usepackage{caption}
\usepackage{subcaption}
%\usepackage{parskip}
\usepackage[skip=12pt plus1pt]{parskip}
\usepackage{pdfpages} % Incluir PDF en documento en LATEX
\usepackage{verbatim} % comentarios
\usepackage{algpseudocode}
\usepackage{algorithm}
\usepackage{pdflscape}
\usepackage{multirow}
\usepackage{afterpage}
\usepackage{array,booktabs,ragged2e}
%\newcolumntype{R}[1]{>{\RaggedLeft\arraybackslash}p{#1}}
\newcolumntype{L}[1]{>{\raggedright\let\newline\\\arraybackslash\hspace{0pt}}m{#1}}
\newcolumntype{C}[1]{>{\centering\let\newline\\\arraybackslash\hspace{0pt}}m{#1}}
\newcolumntype{R}[1]{>{\raggedleft\let\newline\\\arraybackslash\hspace{0pt}}m{#1}}

\floatname{algorithm}{Algoritmo}
\renewcommand{\listalgorithmname}{Lista de algoritmos}
\renewcommand{\algorithmicrequire}{\textbf{Entrada:}}
\renewcommand{\algorithmicensure}{\textbf{Salida:}}

% Comentario con respecto a las referencias:
% Por defecto este documento utiliza el formato IEEEtr para formatear las referencias. Si algun asesor requiere formato APA, solo comente la siguiente linea y descomente la linea debajo. Tambien para que las citas funcionen de manera adecuada, el paquete label debe tener las opciones mostradas, o en su defecto, el año en el listado no lo muestra correctamente \usepackage[english, american]{babel}

% Para utilizar el formato de citas IEEE y comentar los dos parrafos siguientes
\usepackage[backend=bibtex,sorting=none]{biblatex}
% Para utilizar el formato APA, sugiero comentar la linea anterior y descomentar las dos proximas lineas
%\usepackage[backend=biber,style=apa]{biblatex}
%\DeclareLanguageMapping{english}{american-apa}

\makeatletter
\DefineBibliographyExtras{spanish}{%
  \setcounter{smartand}{1}%
  \let\lbx@finalnamedelim=\lbx@es@smartand
  \let\lbx@finallistdelim=\lbx@es@smartand
}
\renewbibmacro*{name:delim:apa:family-given}[1]{%
  \ifnumgreater{\value{listcount}}{\value{liststart}}
    {\ifboolexpr{
       test {\ifnumless{\value{listcount}}{\value{liststop}}}
       or
       test \ifmorenames
     }
       {\printdelim{multinamedelim}}
       {\lbx@finalnamedelim{#1}}}
    {}}
\makeatother


% Estas lineas permiten romper los hipervinculos muy largos en las referencias!!!!
\setcounter{biburllcpenalty}{7000}
\setcounter{biburlucpenalty}{8000}
\addbibresource{x.bib} % ARCHIVO DE BIBLIOGRAFÍA


%\usepackage{url}
\usepackage[bookmarks=true,breaklinks=true,bookmarksopen=false,colorlinks=true,linkcolor=blue]{hyperref}
\usepackage[hyphenbreaks]{breakurl}
% Regla que define explicitamente que caracteres rompen los hipervinculos para separar las lineas
%https://es.overleaf.com/11089898rhgykrqyqytx
% Actualiza en automático la fecha de las citas de internet a la fecha de la compilación del documento
\usepackage{datetime}
\newdateformat{specialdate}{\twodigit{\THEDAY}-\twodigit{\THEMONTH}-\THEYEAR}
%\newdateformat{specialdate}{\twodigit{\THEDAY}-\THEYEAR}
\date{\specialdate\today}

\newcommand{\HRule}{\rule{\linewidth}{0.25mm}}


% CONSTANTES NECESARIAS PARA EL DOCUMENTO ---> MODIFIQUEN A SU CRITERIO
\newcommand{\ncarrera}            {Ingeniería en Tecnologías de la Información}
\newcommand{\nasesorinstitucional}{Dr. SAID POLANCO MARTAGÓN}
\newcommand{\NombreAlumno}{Karina Lizeth Arévalo Flores \\ Carlos Antonio Bocanegra González Dafne Patricia Moreno Flores
Alejandra Carolina Rodríguez Porras \\
Jesús David Treviño Gandarilla}
%Hombres cambien LA por EL
\newcommand{\elolaNombreAlumno}{el}  
\newcommand{\OA}           {o}  %Hombres cambien A por O
\newcommand{\Matricula}           {1730505}  %SU MATRICULA
\newcommand{\NombreProyecto}{Sistema Inteligente de Clasificación de Semillas}
\newcommand{\fechacarta}{2 de octubre de 2024}
\newcommand{\ncuatrimestre}{Septiembre-Diciembre 2024}
\newcommand{\nevalador}{Dr. SAID POLANCO MARTAGÓN}
\newcommand{\FechaExposicion}{2 de octubre de 2024}
\newcommand{\HoraExposionFormatoVenticuatroHoras}{10:00}
\newcommand{\elolaNombreEmpresa}{la}  %Si la empresa es femenino (por ejemplo universidad, usen la) o masculino (el instituto) pongan el
\newcommand{\organismoreceptor}   {Universidad Politécnica de Victoria}

% NOTA: Dos diagonales juntas (\\) indican un saldo de linea. En este caso particular hay 2 (el titulo se ajusta a tres lineas, porque es muy largo. Hacer las adecuaciones pertinentes
\newcommand{\NombreProyectoheader}     {CLASIFICACIÓN DE SEMILLAS.}
%\newcommand{\nasesorempresaria}   {MSI. José Fidencio López Luna}
\newcommand{\fechaPortada}               {Octubre de 2024}





\newcommand{\separacionCorta}{0.0cm}
\newcommand{\separacionLarga}{0.5cm}

\usepackage[overload]{textcase}
\newcommand{\iemph}[1]{\MakeTextUppercase{#1}}

\pagestyle{fancy}
\headheight 45pt
\fancyhead{} % Clear all header fields
\fancyhead[L]{\includegraphics[height=1.00cm]{UTYP.png}}%
\fancyhead[C]{\begin{center}\NombreProyectoheader\end{center}}%
%\fancyhead[R]{\includegraphics[height=1.5cm]{LogoUPV_2019.png}}%
\fancyhead[R]{\includegraphics[height=1.25cm]{LogoUPV_2023.png}}%
\fancyfoot[R]{\thepage} % Clear all footer fields 
\fancyfoot[C]{}
\fancyfoot[L]{}

\DefineBibliographyStrings{english}{%
  references = {Referencias},% replace "references" with "bibliography"  for `book`/`report`
}

\addto\captionsenglish{%
  \renewcommand{\figurename}{Figura}%
  \renewcommand{\tablename}{Tabla}%
} 

\usepackage{wallpaper}
 
 
%\renewcommand{\figurename}{Figura}
%\renewcommand{\tablename}{Tabla}

 
\begin{document}

%-----------------------------------------------------------------------------------------------------------------
% PAGINA 1 - PORTADA
\setcounter{page}{1}
\pagenumbering{roman}
\thispagestyle{empty}

\begin{center}

\begin{tabular}{cp{5cm}c}
\includegraphics[height=2.25cm]{UTYP.png} & 
& \includegraphics[height=2.25cm]{LogoUPV_2023.png}   \\
\end{tabular}

\Large \textbf{UNIVERSIDAD POLITÉCNICA DE VICTORIA}
\vspace{0.5cm}
\hrule
\vspace{0.1cm} 
\hrule
\vspace{0.5cm}


%\HRule \\[\separacionCorta]
\textbf{\iemph{\NombreProyecto}} \\[\separacionLarga]
%\Large \textbf{TESINA}
%\HRule \\[\separacionLarga]
PROYECTO FINAL \\
DE LA CARRERA EN \\
\textbf{\iemph{\ncarrera}} \\[\separacionLarga]

PRESENTA: \\[\separacionCorta]
%\textbf{\Capitalize{\NombreAlumno}\\[\separacionLarga]
\textbf{\iemph{\NombreAlumno}}\\[\separacionLarga]
%EN CUMPLIMIENTO DE \\[\separacionCorta]
%LA ESTADÍA DE LA CARRERA DE \\[\separacionCorta]


DIRECTOR \\[\separacionCorta]
\textbf{\iemph{\nasesorinstitucional}} \\[\separacionCorta]


 MATERIA \\[\separacionCorta]
%\textbf{\iemph{\nasesorempresaria}} \\[\separacionCorta]

ORGANISMO RECEPTOR \\[\separacionCorta]
\textbf{\iemph{\organismoreceptor}} \\[\separacionLarga]

\end{center}
\begin{flushright}
\iemph{Ciudad Victoria, Tamaulipas, \fechaPortada}
\end{flushright}

\HRule 






%-----------------------------------------------------------------------------------------------------------------

%-----------------------------------------------------------------------------------------------------------------

%-----------------------------------------------------------------------------------------------------------------


%-----------------------------------------------------------------------------------------------------------------


%-----------------------------------------------------------------------------------------------------------------


%-----------------------------------------------------------------------------------------------------------------
% PAGINA 7 - INDICE

\clearpage
\addcontentsline{toc}{section}{Índice}
\renewcommand\contentsname{Índice}
\tableofcontents

%-----------------------------------------------------------------------------------------------------------------
% CAPITULOS


\clearpage
\pagenumbering{arabic}
\setcounter{page}{1}
%\input{Capitulo1.tex}

\clearpage
%\input{Capitulo2.tex}

\clearpage
\section{Introducción}
La clasificación de semillas ha evolucionado significativamente a lo largo del tiempo. En sus inicios, los agricultores dependían exclusivamente de la inspección visual y manual, un proceso lento y propenso a errores. A mediados del siglo XX, se introdujeron métodos mecánicos como tamices y separadores por peso, que mejoraron la eficiencia pero aún carecían de precisión para características sutiles. En la década de 1980, surgieron las primeras aplicaciones de visión por computadora en la agricultura, marcando el inicio de una nueva era. Sin embargo, estos sistemas eran costosos y complejos, limitando su adopción generalizada. 
En el último lustro, la visión por ordenador ha sido ampliamente utilizada en el área de la ingeniería agromótica para la inspección, seguimiento, riego y cosecha de los productos agrícolas. Actualmente se están cambiando las operaciones manuales tradicionales como son la mano de obra lenta y propensa a errores de los seres humanos \cite{hernandez2016comparacion}. La clasificación precisa de semillas tiene un impacto económico sustancial en la industria agrícola. Según un estudio realizado por la FAO en 2020, las pérdidas debidas a la clasificación incorrecta de semillas pueden alcanzar hasta el 10\% de la producción total en algunos cultivos \cite{mensah2021adoption}. Esto se traduce en millones de dólares en pérdidas anuales para los agricultores y la industria alimentaria. Por otro lado, la implementación de métodos de clasificación más precisos ha demostrado aumentar la productividad entre un 15\% y un 25\%, dependiendo del cultivo. Por ejemplo, en un estudio de caso realizado en campos de trigo en Estados Unidos, la adopción de tecnologías avanzadas de clasificación de semillas resultó en un aumento del rendimiento del 18\% y una reducción de costos del 12\% en un período de tres años. 



La clasificación precisa de semillas es fundamental para garantizar la calidad y el rendimiento de los cultivos. Este proyecto se enfoca en desarrollar un sistema de clasificación de semillas utilizando una cámara ESP32-CAM, que combina tecnología de visión por computadora con un dispositivo compacto y económico. La motivación principal de esta investigación es mejorar la eficiencia y precisión en la selección de semillas, lo que puede tener un impacto significativo en la productividad agrícola. Uno de los principales retos que se enfrentarán es la gran variabilidad en las características físicas de las semillas, incluyendo tamaño, color y forma. Estas variaciones pueden ser sutiles entre diferentes variedades de la misma especie, o incluso entre semillas de la misma variedad debido a factores ambientales durante su desarrollo. Además, las condiciones de iluminación juegan un papel crucial en la captura de imágenes precisas. La iluminación inconsistente o inadecuada puede alterar la percepción del color y la textura de las semillas, lo que podría llevar a clasificaciones erróneas. Otro desafío importante es la necesidad de desarrollar algoritmos robustos capaces de manejar estas variaciones y adaptarse a diferentes tipos de semillas sin requerir una recalibración constante. Estos algoritmos deben ser lo suficientemente flexibles para distinguir entre características críticas y variaciones normales, y al mismo tiempo ser eficientes para procesar grandes volúmenes de semillas en tiempo real. La integración de estas capacidades en un dispositivo compacto y económico como la ESP32-CAM añade una capa adicional de complejidad, requiriendo una optimización cuidadosa de los recursos computacionales y de memoria.

Este proyecto se alinea estrechamente con los objetivos de agricultura sostenible y seguridad alimentaria establecidos por las Naciones Unidas en sus Objetivos de Desarrollo Sostenible (ODS). La clasificación precisa de semillas contribuye directamente al ODS 2 (Hambre Cero) y al ODS 12 (Producción y Consumo Responsables). Al mejorar la eficiencia en la selección de semillas, se reduce el desperdicio de recursos como agua, tierra y fertilizantes, lo que a su vez disminuye la huella ecológica de la producción agrícola. Además, la mejora en la calidad de las semillas seleccionadas puede llevar a cultivos más resistentes a enfermedades y condiciones climáticas adversas, contribuyendo así a la resiliencia alimentaria frente al cambio climático. La FAO estima que la adopción de tecnologías como la propuesta en este proyecto podría contribuir a reducir el hambre en un 20\% en regiones en desarrollo para el año 2030, al tiempo que disminuye el uso de agua en la agricultura en un 15\% \cite{trendov2019tecnologias}.

En el campo de la visión por computadora y el aprendizaje automático aplicado a la agricultura, existen diversos algoritmos y librerías que pueden ser fundamentales para el desarrollo de un sistema de clasificación de semillas utilizando una ESP32-CAM. Estos componentes software son esenciales para procesar las imágenes capturadas y realizar una clasificación precisa y eficiente. Entre los algoritmos más relevantes para este tipo de proyecto se encuentran los de procesamiento de imágenes y los de aprendizaje automático. En cuanto a los algoritmos de clasificación, las Máquinas de Soporte Vectorial (SVM) y los árboles de decisión son opciones populares debido a su eficacia y relativa ligereza computacional. Para implementar estos algoritmos, varias librerías de código abierto son invaluables. OpenCV, una de las más utilizadas en visión por computadora, ofrece una amplia gama de funciones para el procesamiento de imágenes y la implementación de algoritmos de visión. Su versión para sistemas embebidos, OpenCV-lite, es especialmente relevante para dispositivos con recursos limitados como la ESP32-CAM \cite{garcia2015learning}.
TensorFlow Lite es otra librería crucial, diseñada específicamente para ejecutar modelos de aprendizaje automático en dispositivos con restricciones de recursos. Esta librería permite la implementación eficiente de redes neuronales y otros algoritmos de aprendizaje automático en la ESP32-CAM.
Para el desarrollo del software en la ESP32-CAM, el framework ESP-IDF (Espressif IoT Development Framework) proporciona un conjunto de herramientas y librerías optimizadas para este hardware específico. Este framework facilita la integración de los algoritmos de visión por computadora con las capacidades de la cámara y el procesamiento de la ESP32-CAM. Además, librerías como NumPy y SciPy, aunque no se ejecutan directamente en la ESP32-CAM, son invaluables para el desarrollo y entrenamiento de modelos en un entorno de desktop antes de su implementación en el dispositivo embebido. Estas librerías ofrecen funciones avanzadas de cálculo numérico y procesamiento de señales que son fundamentales en el desarrollo de algoritmos de clasificación.

La selección y adaptación adecuada de estos algoritmos y librerías es crucial para desarrollar un sistema de clasificación de semillas eficiente y preciso utilizando la ESP32-CAM. El desafío radica en optimizar estos componentes software para que funcionen de manera efectiva dentro de las limitaciones de memoria y potencia de procesamiento del dispositivo, manteniendo al mismo tiempo un alto nivel de precisión en la clasificación.

El objetivo principal de este estudio es diseñar y implementar un sistema automatizado capaz de clasificar diferentes tipos de semillas con alta precisión y contabilizarlas. Además, se busca crear una solución accesible y fácil de usar para agricultores y productores de semillas, que pueda integrarse fácilmente en sus procesos actuales. Este enfoque tiene el potencial de reducir el tiempo y los recursos necesarios para la clasificación manual de semillas, al tiempo que minimiza los errores humanos.
El problema que se pretende resolver es la dificultad y el costo asociados con la clasificación manual de semillas, que a menudo resulta en inexactitudes y requiere mucho tiempo. Los métodos tradicionales de clasificación pueden ser subjetivos y propensos a errores, lo que afecta la calidad de los cultivos y, en última instancia, la producción agrícola. Al automatizar este proceso mediante el uso de tecnología de visión por computadora, se espera mejorar significativamente la precisión y la eficiencia de la clasificación de semillas.


Este proyecto se basa en investigaciones previas en el campo de la visión por computadora aplicada a la agricultura. Por ejemplo, el trabajo de García-Mateos \cite{garcia2015study} demostró la eficacia de los sistemas de visión por computadora para la clasificación de frutas y verduras. Además, estudios como el de Mahajan \cite{mahajan2015image} han explorado el uso de técnicas de aprendizaje automático para la clasificación de semillas, sentando las bases para nuestra investigación.
Al combinar estos avances con la versatilidad y accesibilidad de la ESP32-CAM, este proyecto busca crear una solución innovadora que pueda tener un impacto real en la industria agrícola. Se espera que los resultados de esta investigación no solo mejoren la eficiencia en la clasificación de semillas, sino que también abran nuevas posibilidades para la aplicación de tecnologías de visión por computadora en otros aspectos de la agricultura.



\clearpage
\section{Antecedentes teóricos}

En las últimas décadas, la aplicación de tecnologías de visión por computadora para la clasificación y detección de semillas ha experimentado un desarrollo significativo. Diversos investigadores han abordado este desafío utilizando diferentes enfoques y tecnologías. La implementación de sistemas automatizados para la clasificación de semillas representa un avance importante en la agricultura de precisión, mejorando la eficiencia y precisión en los procesos de selección. Uno de los trabajos pioneros en este campo fue realizado por Kumar y Singh \cite{qiu2018variety}, quienes desarrollaron un sistema de clasificación de semillas utilizando técnicas de procesamiento de imágenes y redes neuronales convolucionales. Su investigación demostró una precisión del 95\% en la clasificación de diferentes tipos de semillas, estableciendo una base sólida para futuros desarrollos en esta área. En el contexto de sistemas embebidos, Patil \cite{patil2020realtime} implementó un sistema de bajo costo utilizando una ESP32-CAM para la clasificación de frutas en tiempo real. Su trabajo demostró la viabilidad de utilizar dispositivos de recursos limitados para tareas de visión por computadora, alcanzando una precisión del 89\% en condiciones controladas de iluminación. La clasificación automatizada de semillas ha sido abordada desde diferentes enfoques. Por ejemplo, Zhang et al. (2021) utilizaron técnicas de aprendizaje profundo para clasificar semillas de arroz, logrando una precisión del 98\%. Su sistema empleaba características como forma, tamaño y textura para la clasificación, aunque requería hardware más potente que una ESP32-CAM \cite{jin2022identification}. 

\subsubsection{Antecedentes históricos de la clasificación de semillas}
A lo largo de la historia, la clasificación de semillas ha sido un proceso crucial en la agricultura, ya que influye directamente en la calidad de los cultivos y, por ende, en el rendimiento de las cosechas. Originalmente, este trabajo se realizaba de manera completamente manual. Los agricultores, mediante una inspección visual, eran los encargados de separar las semillas basándose en características visibles como el tamaño, el color o la forma. Sin embargo, este método, además de ser laborioso y tedioso, estaba sujeto a errores debido a la variabilidad en la percepción humana. Lo que para una persona podía parecer una semilla adecuada, para otra podría no serlo, y esa falta de uniformidad en la selección resultaba en una calidad inconsistente en las cosechas.

Con el tiempo y el avance de la mecanización agrícola a mediados del siglo XX, comenzaron a implementar herramientas y dispositivos que facilitan esta labor. Uno de los primeros avances fue el uso de tamices, herramientas parecidas a coladores grandes que permitían separar las semillas según su tamaño. También se introdujeron separadores por peso, los cuales identifican y clasifican las semillas dependiendo de su densidad. Estos dispositivos representan una mejora significativa en la eficiencia del proceso, pero aún presentaban limitaciones. Por ejemplo, no podían diferenciar con precisión entre semillas que, aunque tenían el mismo tamaño o peso, variaba en características más sutiles, como la forma o el color, aspectos que también influyen en la calidad de la semilla.

El verdadero cambio en la clasificación de semillas ocurrió con el desarrollo de tecnologías más avanzadas en las décadas posteriores. En los años 80, la llegada de la visión por computadora marcó un punto de inflexión. Esta tecnología permitió que las máquinas "vieran" las semillas de una manera mucho más precisa, utilizando cámaras y sistemas computacionales capaces de analizar imágenes y detectar diferencias que serían imperceptibles para el ojo humano. Si bien en sus inicios estos sistemas son costosos y difíciles de implementar, sentaron las bases para las futuras innovaciones en el campo de la agricultura.

Con el tiempo, la integración de la inteligencia artificial (IA) y el aprendizaje automático (machine learning) ha mejorado aún más los sistemas de clasificación. A través de estos avances, las máquinas no solo pueden ver las semillas, sino que también pueden aprender de sus características y realizar clasificaciones más rápidas y precisas. Estos sistemas automatizados son capaces de analizar en tiempo real atributos como el tamaño, la forma y el color, permitiendo una inspección detallada que sería imposible de realizar de forma manual con la misma velocidad y precisión.

Hoy en día, la visión por computadora y la IA se han combinado para ofrecer soluciones que permiten una mayor precisión en la clasificación de semillas, lo que reduce significativamente los errores humanos y aumenta la eficiencia en la agricultura. La integración de estas tecnologías ha mejorado enormemente la capacidad de los agricultores para seleccionar las mejores semillas y, por lo tanto, obtener mejores cosechas.

\subsubsection{El uso de dispositivos embebidos y la democratización de la tecnología}
Más recientemente, el uso de dispositivos embebidos como la ESP32-CAM ha permitido llevar estas innovaciones tecnológicas a un nivel más accesible para los agricultores. La ESP32-CAM, un pequeño y económico dispositivo que incluye una cámara y conectividad Wi-Fi, ha demostrado ser una herramienta valiosa para la clasificación de semillas en tiempo real. A diferencia de los sistemas de visión por computadora más antiguos y costosos, la ESP32-CAM permite realizar análisis de imágenes con una inversión mucho menor, lo que ha democratizado el acceso a esta tecnología.

Este dispositivo, cuando se combina con bibliotecas de software como OpenCV y TensorFlow Lite, es capaz de ejecutar algoritmos de procesamiento de imágenes y aprendizaje automático, incluso en entornos de recursos limitados. Gracias a estos avances, la ESP32-CAM puede tomar fotos de las semillas, procesarlas y clasificarlas casi de manera instantánea, permitiendo que los agricultores trabajen con una mayor eficiencia sin necesidad de equipos complejos y costosos.

El impacto de esta tecnología en la agricultura moderna es significativo. Permite a los agricultores no solo ahorrar tiempo, sino también optimizar el uso de recursos, lo que se traduce en una mayor sostenibilidad. Además, al reducir la dependencia de la mano de obra para tareas repetitivas y tediosas, los agricultores pueden centrar sus esfuerzos en otras áreas importantes de la producción.

\subsubsection{Proyectos similares}
En los últimos años, se han desarrollado diversos sistemas de visión artificial utilizando dispositivos embebidos como la ESP32-CAM para tareas de clasificación y conteo. A continuación, se presentan algunos trabajos relevantes en este campo:

\begin{itemize}
    \item \textbf{Clasificación de Frutas:} Kim \cite{kim2022embedded} desarrollaron un sistema de clasificación de frutas utilizando una ESP32-CAM y redes neuronales convoluciones ligeras. Su sistema logró una precisión del 94.3\% en la clasificación de 5 tipos diferentes de frutas, demostrando la viabilidad de implementar modelos de deep learning en dispositivos con recursos limitados.
    \item \textbf{Conteo Automatizado de Productos:} Martínez-López \cite{martinez2023automated} implementaron un sistema de conteo automatizado para una línea de producción utilizando una ESP32-CAM. El sistema procesaba las imágenes en tiempo real y enviaba los datos a una base de datos en la nube, alcanzando una precisión del 98\% en el conteo de objetos en movimiento.

    \item \textbf{Clasificación de Granos:} Un trabajo particularmente relevante es el de Zhang \cite{zhang2021rice}, quienes desarrollaron un sistema de clasificación de granos de arroz utilizando una ESP32-CAM. Su enfoque combinó técnicas de procesamiento de imágenes tradicionales con un modelo de machine learning ligero, logrando una precisión del 91\% en la clasificación de diferentes variedades de arroz.
     
    \item \textbf{Control de Calidad en Agricultura:} Patel y Singh \cite{patel2023quality} implementaron un sistema de control de calidad para la clasificación de tomates utilizando una ESP32-CAM. Su sistema podía detectar defectos superficiales y clasificar los tomates según su madurez con una precisión del 89\%.
\end{itemize}

\clearpage
\section{Descripción del problema}
La clasificación de semillas, uno de los pasos más críticos en la agricultura, afecta directamente a la calidad de la cosecha. Hasta ahora, la selección de semillas en función de características visibles como el tamaño, la forma, el color y la textura ha sido su naturaleza. La mayoría de las veces, también se realiza de forma manual. Sin embargo, la clasificación manual es muy imprecisa, hace perder tiempo y es costosa. Debido a las limitaciones de la precisión en el tiempo y el costo de la clasificación manual, los agricultores experimentan enormes dificultades debido a la disparidad física de las semillas, la fatiga visual y la subjetividad de la percepción humana. Esta clasificación no uniforme da lugar a enormes errores, por lo tanto, en la calidad de la cosecha producida y su producción agrícola final.Los errores de clasificación pueden provocar una menor germinación, menores rendimientos de los cultivos, mayor vulnerabilidad a las enfermedades y, en general, una calidad inferior en los productos agrícolas.
\\
Además, la población mundial sigue aumentando, lo que aumenta la demanda mundial de alimentos y la necesidad de eficiencia y precisión en los métodos de producción. La necesidad de una clasificación precisa de las semillas maximiza el rendimiento de los cultivos y minimiza el desperdicio. Si bien ha habido avances tecnológicos en la maquinaria agrícola en los últimos tiempos, existen varios países en desarrollo donde los agricultores aún realizan el proceso de forma manual o utilizando equipos muy simples diseñados para ese fin. Esto se debe a que los equipos de clasificación de semillas automatizados son muy caros y complicados, están diseñados principalmente para grandes explotaciones y están fuera de su alcance. Por esta razón, existen muchas ineficiencias en la producción agrícola con pérdidas económicas bastante grandes.
\\Uno de los desafíos clave al desarrollar un sistema de clasificación automatizado de semillas son las diferencias extremadamente altas en las características físicas, incluso dentro de la misma especie. Es decir, existen factores muy variados como el tamaño, la forma, y color, que tienen su influencia significativa. Su variabilidad se debe tanto a factores genéticos como ambientales. Además, las condiciones de captura fotográfica como la iluminación o el ángulo de la toma también son críticas. Por lo tanto, una característica esencial de un sistema eficaz es su capacidad de tomar en cuenta todos estos factores y proporcionar resultados consistentes en varias condiciones.
Por último,  está el costo de la implementación. Los agricultores de pequeños y medianos campos no pueden financiar la inversión en tecnologías avanzadas, lo que mantiene el uso de la clasificación manual. En esta situación, el sistema de clasificación de semillas basado en tecnología accesible y económica, como ESP32, tiene sentido. Dado que el dispositivo puede compartir la funcionalidad de la grabación de imágenes con un microcontrolador económico pero eficiente, ofrece una oportunidad para la implementación de la clasificación automatizada a un costo mucho menor al de los sistemas tradicionales.

\clearpage
\section{Motivación}

La necesidad de mejorar la clasificación de semillas, es un proceso fundamental para aumentar la productividad agrícola, es la principal motivación de este proyecto. El uso de tecnologías emergentes como el aprendizaje automático ha demostrado ser efectivo para optimizar los costos y aumentar la producción, lo que impulsa el uso de estas herramientas en la agricultura. Se busca crear un sistema accesible para los pequeños agricultores que les permita automatizar tareas que tradicionalmente requerían un gran esfuerzo manual. Se utilizará un dispositivo compacto y de bajo costo como la ESP32-CAM.

La promoción de una agricultura más sostenible es otra motivación importante. Al mejorar la clasificación de semillas, se logra un uso más eficiente de recursos vitales como el agua, la tierra y los fertilizantes, lo que contribuye a una producción más responsable con el medio ambiente. Además, este proyecto podría abrir la puerta a nuevas aplicaciones de la visión por computadora en otros aspectos de la agricultura, como la recolección o el monitoreo de plagas, extendiendo los beneficios tecnológicos a diferentes áreas del campo.

Finalmente, este proyecto se alinea con la agenda 2030 de los Objetivos de Desarrollo Sostenible, particularmente con el ODS 2 (Hambre Cero) y el ODS 12 (Producción y Consumo Responsables). La promoción de una agricultura más sostenible es otra motivación importante. Al mejorar la clasificación de semillas, se logra un uso más eficiente de recursos vitales como el agua, la tierra y los fertilizantes, lo que contribuye a una producción más responsable con el medio ambiente. Además, este proyecto podría abrir la puerta a nuevas aplicaciones de la visión por computadora en otros aspectos de la agricultura, como la recolección o el monitoreo de plagas, extendiendo los beneficios tecnológicos a diferentes áreas del campo.


\clearpage
\section{Justificación}

En la agricultura moderna, la clasificación precisa de semillas es una necesidad de alta precisión, ya que tiene un impacto directo en la calidad de los cultivos y en la eficiencia de la producción agrícola. Actualmente, la mayoría de los métodos de clasificación de semillas son manuales o mecánicos, lo que los hace propensos a errores humanos y con baja eficiencia en la diferenciación de características sutiles. La productividad agrícola y, en consecuencia, el ingreso de los agricultores, se ven afectados negativamente por esta situación. La FAO ha informado que, en algunos cultivos, la clasificación incorrecta puede generar pérdidas que alcanzan hasta el 10 \% de la producción total. Implementar el sistema automatizado propuesto en este proyecto, que utiliza visión por computadora, no solo mejoraría la precisión en la clasificación, sino que también reduciría costos al optimizar la utilización de los recursos en la producción.


El objetivo de este proyecto es automatizar el proceso de clasificación de semillas, ya que esto tiene un impacto directo en la productividad agrícola, la reducción de costos y la sostenibilidad ambiental. La combinación de tecnologías avanzadas como la visión por computadora y el aprendizaje automático en un dispositivo compacto y económico como ESP32-CAM ofrece una solución innovadora y accesible para los agricultores, especialmente en regiones donde el acceso a tecnología de alto costo es limitado. Además, al disminuir el trabajo manual, los agricultores obtienen una mejor calidad de vida y pueden concentrarse en otras actividades productivas. Los pequeños agricultores obtienen mayores beneficios al reducir los costos de clasificación de semillas, lo que fortalece las economías locales y mejora la seguridad alimentaria en sus comunidades.



\clearpage
\section{Alcances y limitaciones}

El presente proyecto se compromete a desarrollar un prototipo funcional para la clasificación automatizada de semillas utilizando una cámara ESP32-CAM, con el objetivo de mejorar la eficiencia en la selección de semillas. El alcance del trabajo incluye la implementación de un sistema capaz de capturar imágenes de las semillas, procesarlas mediante algoritmos de visión por computadora y clasificarlas de acuerdo con parámetros predefinidos como su tamaño, color y forma. El sistema será probado en condiciones controladas de laboratorio para garantizar su funcionamiento dentro de un entorno predecible y adecuado. Además, se espera que esta solución sea accesible económicamente, permitiendo que pequeños agricultores puedan beneficiarse de una tecnología que tradicionalmente ha sido de alto costo.

Sin embargo, el proyecto presenta varias limitaciones que deben ser consideradas. La primera de ellas está relacionada con las condiciones de iluminación. Dado que el sistema depende de la captura de imágenes para su funcionamiento, cualquier variación en la iluminación podría afectar negativamente la calidad de las imágenes obtenidas y, por lo tanto, la precisión de la clasificación. El control de la iluminación es limitado a las condiciones de prueba en un entorno específico, y el rendimiento del sistema en ambientes no controlados podría no ser tan fiable.

Otra limitación importante es la capacidad del hardware. La ESP32-CAM es un dispositivo compacto y económico, pero sus recursos de procesamiento y memoria son limitados. Esto restringe el tipo y la complejidad de los algoritmos de aprendizaje automático y visión por computadora que se pueden implementar. Aunque se buscará optimizar el uso de recursos, las limitaciones inherentes del hardware podrían afectar la velocidad y precisión del sistema, especialmente al manejar grandes volúmenes de datos o imágenes de alta resolución.

Además, la variabilidad en las características físicas de las semillas representa otro reto. Las semillas pueden variar considerablemente en términos de tamaño, color y textura, incluso dentro de la misma especie. Este nivel de variabilidad puede dificultar que el sistema logre una clasificación precisa sin ajustes o recalibraciones adicionales, lo que podría limitar su aplicación en ciertos tipos de semillas o en condiciones más extremas.

Finalmente, el proyecto también está condicionado por los recursos materiales y de tiempo disponibles. El uso de un dispositivo de bajo costo implica que no se podrá contar con tecnologías adicionales que podrían mejorar la precisión del sistema, como sensores complementarios o algoritmos más avanzados. Además, las pruebas estarán restringidas a un entorno y periodo específico, lo que limita la capacidad de evaluar el sistema en diferentes contextos geográficos o ambientales. Estos factores podrían restringir el alcance del proyecto y generar la necesidad de futuras investigaciones para adaptar y mejorar el sistema a condiciones más diversas.





%-----------------------------------------------------------------------------------------------------------------
% REFERENCIAS


\clearpage
%Let's cite! The Einstein's journal paper \cite{dirac} and the Dirac's 
%book \cite{einstein} are physics related items. 

%\Urlmuskip=0mu plus 1mu\relax
\addcontentsline{toc}{section}{Referencias} 
\printbibliography
 
\end{document}
